\inicializarSubseccion{Conclusión de Viviana}

\tab Mi mayor aprendizaje en este proyecto fue ver cómo se aplicaban las teorías de conjuntos a la hora de hacer las preguntas de nuestro formulario y como en base a las respuestas las podiamos analizar con la teoría de conjuntos, ya que antes de esto cuando tenía que hacer encuestas o formularios nunca lo hacía en base a una estructura y poniendo preguntas y respuestas clave como seleccionar mínimo más de una o selecciona todas las que aplican etc, si no que simplemente lo ponía cuando yo sentía que aplicaba o si me acordaba pero nunca pensando si realmente esas opciones aplican para el tipo de pregunta. Por el lado de código me gustó mucho y aprendí bastante al hacer el florero y sus respectivas funciones par poder hacer el modelo en 2d y 3d ya que con esta base me facilitó mucho saber cómo aplicarlo a nuestro proyecto ya que en este batallé al momento de hacer las funciones para el volumen ya que en un principio mi volumen me salí negativo pero no fue si no que hasta que me puse a leer las instrucciones detalladamente y logré darme cuenta que el ejemplo que usaban era una ecuación de grado 3 y la mía era de grado 4  y grado 2 entonces corregí las ecuaciones necesarias agregando o quitando elevados y listo.Con los cálculos pudimos determinar el material ya que un principio teníamos un costó muy elevado por el material que teníamos fuera de eso no los llegamos a necesitar tanto.
